\section*{الفصل \ref{ch:g10:why-clouds-form} --- لماذا تتكون السحب}

\begin{solution}{exo:g10:why-clouds-form:1}
\omterm{def:g10:why-clouds-form:humidity}{الرطوبة النسبية} هي محتوى الهواء من بخار الماء مقسومًا على محتواه عند \omterm{def:g10:why-clouds-form:humidity}{التشبع} في
درجة الحرارة نفسها والضغط نفسه، معبَّرًا عنها بنسبة مئوية. \omterm{def:g10:why-clouds-form:dew-point}{ونقطة الندى} هي درجة
الحرارة التي يجب أن يُبرَّد إليها الهواء، مع المحافظة على محتواه من الماء وعلى
ضغطه، كي يصير \omterm{def:g10:why-clouds-form:humidity}{مشبعًا}.
\end{solution}

\begin{solution}{exo:g10:why-clouds-form:2}
\num{3.8} و\num{7.6} و\num{14.7} و\qty{27.2}{g/kg}. وكل احترار بمقدار
\qty{10}{\degreeCelsius} يضربه في نحو 2 ($7.6/3.8 = 2.0$،
$14.7/7.6 \approx 1.9$، $27.2/14.7 \approx 1.9$).
\end{solution}

\begin{solution}{exo:g10:why-clouds-form:3}
المحتوى عند \omterm{def:g10:why-clouds-form:humidity}{التشبع} عند \qty{24}{\degreeCelsius}: \qty{18.9}{g/kg}؛ \omterm{def:g10:why-clouds-form:humidity}{والرطوبة
النسبية} $10.0/18.9 \approx \qty{53}{\percent}$. ويكون المحتوى عند \omterm{def:g10:why-clouds-form:humidity}{التشبع}
\qty{10.0}{g/kg} عند \qty{14}{\degreeCelsius}: فتلك هي \omterm{def:g10:why-clouds-form:dew-point}{نقطة الندى}.
\end{solution}

\begin{solution}{exo:g10:why-clouds-form:4}
في الصيف يكون الهواء رطبًا \omterm{def:g10:why-clouds-form:dew-point}{ونقطة نداه} أعلى بكثير من درجة حرارة الكأس: فالهواء
الملامس للكأس يبرد إلى ما دون \omterm{def:g10:why-clouds-form:dew-point}{نقطة نداه} فيتكاثف الماء عليها. أما هواء الشتاء
الجاف جدًا داخل البيت \omterm{def:g10:why-clouds-form:dew-point}{فنقطة نداه} أدنى من درجة حرارة الكأس: فحتى حين يبرد
بملامستها، لا يصير \omterm{def:g10:why-clouds-form:humidity}{مشبعًا}.
\end{solution}

\begin{solution}{exo:g10:why-clouds-form:5}
التسخين من الأسفل (ركام بعد الظهر فوق أرض مشمسة)؛ والجبال (ريح تُجبر على الصعود
على سلسلة جبلية)؛ والجبهات (هواء دافئ يصعد فوق هواء بارد)؛ والتقارب (رياح تلتقي
قرب خط الاستواء أو تتدفق نحو منخفض جوي).
\end{solution}

\begin{solution}{exo:g10:why-clouds-form:6}
$125 \times (30 - 18) = \qty{1500}{m}$؛ $125 \times (22 - 20) =
\qty{250}{m}$. اليوم الثاني: إذ لا تتكون السحب إلا على ارتفاع \qty{250}{m} فوق
الأرض.
\end{solution}

\begin{solution}{exo:g10:why-clouds-form:7}
على ارتفاع \qty{1}{km}: $18 - 9.8 = \qty{8.2}{\degreeCelsius}$؛ وعلى \qty{2}{km}:
$18 - 19.6 = \qty{-1.6}{\degreeCelsius}$. وفي هبوطها تسخن بالمعدل نفسه فتبلغ
الأرض عند \qty{18}{\degreeCelsius} من جديد.
\end{solution}

\begin{solution}{exo:g10:why-clouds-form:8}
مع \qty{4}{\degreeCelsius/km}، تبرد الكتلة غير المشبعة
(\qty{9.8}{\degreeCelsius/km}) والكتلة \omterm{def:g10:why-clouds-form:humidity}{المشبعة}
(\qty{5}{\degreeCelsius/km}) كلتاهما أسرع من محيطهما: \omterm{def:g10:why-clouds-form:stability}{فالهواء مستقر} في
الحالتين. ومع \qty{12}{\degreeCelsius/km}، تبرد حتى الفقاعة غير المشبعة أبطأ من
محيطها: فتبقى أدفأ منه وتواصل الصعود، وتصعد فقاعات الهواء الدافئ بحرية.
\end{solution}

\begin{solution}{exo:g10:why-clouds-form:9}
$14.7 - 7.6 = \qty{7.1}{g}$ لكل كيلوغرام من الهواء. \omterm{def:g10:earth-energy-budget:latent-heat}{والحرارة الكامنة} تدفئ
الهواء الذي يتكاثف فيه الماء، فتبطئ تبرده.
\end{solution}

\begin{solution}{exo:g10:why-clouds-form:10}
تكبر الأحجام كمكعب القطر: $(0.5/0.02)^3 = 25^3 = 15\,625$
قطيرة لقطرة رذاذ؛ و$(2/0.02)^3 = 100^3 = 1\,000\,000$ لقطرة
مطر.
\end{solution}

\begin{solution}{exo:g10:why-clouds-form:11}
تحت سماء صافية تشع الأرض أشعة تحت حمراء فتبرد، ويبرد معها الهواء الملامس لها.
وهذا الهواء البارد أكثف من الهواء الذي فوقه؛ فينحدر على السفوح ويتجمع في قاع
الوادي. وهناك يواصل التبرد طوال الليل حتى يبلغ \omterm{def:g10:why-clouds-form:dew-point}{نقطة نداه}، فتتكون قطيرات في الهواء
نفسه: إنه الضباب. أما على قمم التلال فالهواء ليس بهذه البرودة ويبقى فوق \omterm{def:g10:why-clouds-form:dew-point}{نقطة
نداه}.
\end{solution}

\begin{solution}{exo:g10:why-clouds-form:12}
على الارتفاع $z$ تكون درجة الحرارة $T - a z$ \omterm{def:g10:why-clouds-form:dew-point}{ونقطة الندى} $T_d - b z$؛
وتتساويان حين $(a - b) z = T - T_d$، إذن $z = (T - T_d)/(a - b)$. ومع
$b = 1.7$: $a - b = 8.1$، و$1000/8.1 \approx \qty{123}{m}$ لكل درجة.
\end{solution}

\begin{solution}{exo:g10:why-clouds-form:13}
\omterm{def:g10:why-clouds-form:condensation-level}{مستوى التكاثف}: $125 \times (15 - 11) = \qty{500}{m}$؛ ودرجة الحرارة عنده:
$15 - 9.8 \times 0.5 \approx \qty{10.1}{\degreeCelsius}$. القمة: $10.1 - 5
\times 1.5 = \qty{2.6}{\degreeCelsius}$. ونزولًا \qty{2}{km} بالمعدل الجاف:
$2.6 + 9.8 \times 2 \approx \qty{22.2}{\degreeCelsius}$، أي أدفأ بنحو سبع درجات
مما كان على الجانب المواجه للريح.
\end{solution}

\begin{solution}{exo:g10:why-clouds-form:14}
عند \qty{-20}{\degreeCelsius}: $126/103 \approx 1.22$، أي فوق \omterm{def:g10:why-clouds-form:humidity}{تشبع} بنسبة
\qty{22}{\percent} بالنسبة إلى الجليد؛ وعند \qty{-10}{\degreeCelsius}: $287/260 \approx
1.10$، أي \qty{10}{\percent}. وفي حدود درجات الحرارة هذه، كلما كانت السحابة أبرد
كانت أفضلية الجليد أكبر: فالآلية تعمل على أفضل وجه في الجزء العلوي البارد من
السحابة حيث تختلط القطيرات بالبلورات.
\end{solution}

\begin{solution}{exo:g10:why-clouds-form:15}
غير مشبعة: تبرد الكتلة بمقدار \qty{9.8}{\degreeCelsius}، ومحيطها بمقدار
\qty{6.5}{\degreeCelsius}: فهي أبرد منه بمقدار \qty{3.3}{\degreeCelsius} فتهبط
عائدة. \omterm{def:g10:why-clouds-form:humidity}{مشبعة}: تبرد بمقدار \qty{5}{\degreeCelsius}، فهي أدفأ بمقدار
\qty{1.5}{\degreeCelsius} وتواصل الصعود. \omterm{def:g10:why-clouds-form:stability}{فالهواء غير مستقر} بشرط واحد فقط: أن
يكون الهواء الصاعد \omterm{def:g10:why-clouds-form:humidity}{مشبعًا}.
\end{solution}

\begin{solution}{pb:g10:why-clouds-form:1}
\textbf{1.} \qty{14.7}{g/kg}.

\textbf{2.} $11.4/14.7 \approx \qty{78}{\percent}$.

\textbf{3.} \qty{16}{\degreeCelsius}، حيث يبلغ المحتوى عند \omterm{def:g10:why-clouds-form:humidity}{التشبع}
\qty{11.4}{g/kg}.

\textbf{4.} $11.4/24.1 \approx \qty{47}{\percent}$.

\textbf{5.} لا: فهو يحمل \qty{11.4}{g/kg} نفسها. وقد انخفضت رطوبته النسبية لأنه،
وقد صار أدفأ، صار قادرًا على حمل المزيد.

\textbf{6.} عند \omterm{def:g10:why-clouds-form:dew-point}{نقطة نداه}، \qty{16}{\degreeCelsius}.

\textbf{7.} $11.4 - 8.7 = \qty{2.7}{g}$ لكل كيلوغرام من الهواء.

\textbf{8.} تحت سماء ليلية صافية، يشع العشب وأسطح السيارات أشعة تحت حمراء فتصير
أبرد من الهواء؛ فيبرد الهواء الملامس لها إلى ما دون \omterm{def:g10:why-clouds-form:dew-point}{نقطة نداه} قبل بقية
الهواء.

\textbf{9.} $(28.0 - 13.4)/1.5 \approx \qty{9.7}{\degreeCelsius/km}$،
وهو المعدل الجاف عمليًا: فقد خلطت فقاعات بعد الظهر هذه الطبقة.

\textbf{10.} $(13.4 - 0.4)/2.0 = \qty{6.5}{\degreeCelsius/km}$.

\textbf{11.} مستقر بالنسبة إلى كتلة غير مشبعة، تبرد أسرع
(\qty{9.8}{\degreeCelsius/km})؛ \omterm{def:g10:why-clouds-form:stability}{وغير مستقر} بالنسبة إلى كتلة \omterm{def:g10:why-clouds-form:humidity}{مشبعة}، تبرد
أبطأ (\qty{5}{\degreeCelsius/km}).

\textbf{12.} ترتفع درجة الحرارة مع الارتفاع، من \num{-0.3} إلى
\qty{1.2}{\degreeCelsius}: إنها طبقة مستقرة جدًا تعمل كغطاء. فتتوقف السحابة
الصاعدة هناك: وتقع قممها قرب \qty{4}{km}.

\textbf{13.} $28 - 16 = \qty{12}{\degreeCelsius}$.

\textbf{14.} درجة الحرارة $28 - 9.8 = \qty{18.2}{\degreeCelsius}$؛ \omterm{def:g10:why-clouds-form:dew-point}{ونقطة الندى}
$16 - 1.8 = \qty{14.2}{\degreeCelsius}$.

\textbf{15.} من 12 إلى $18.2 - 14.2 = 4.0$: فقد تقلّص بمقدار
\qty{8}{\degreeCelsius}، أي بمعدل \qty{8}{\degreeCelsius} لكل كيلومتر.

\textbf{16.} يُسدّ فرق درجة واحدة في $1/8$ كيلومتر، أي
\qty{125}{m}.

\textbf{17.} $125 \times (20 - 16) = \qty{500}{m}$.

\textbf{18.} تبقى \omterm{def:g10:why-clouds-form:dew-point}{نقطة الندى} عند \qty{16}{\degreeCelsius} بينما يسخن الهواء قرب
الأرض: فيكبر الفرق، ويكبر معه الارتفاع الذي يُسدّ عنده.

\textbf{19.} $28 - 9.8\,z = 16 - 1.8\,z$.

\textbf{20.} $8\,z = 12$، إذن $z = 1.5$: تقع قاعدة ركام بعد الظهر
\textbf{على نحو \qty{1.5}{km}} فوق السهل (ويعطي الحساب الدقيق بخصائص الماء
\qty{1.51}{km}).
\end{solution}
